\begin{alist}
\item Για να ορίζεται η εξίσωση $(1)$ πρέπει να είναι $x + 1 > 0$ και $1 - x > 0$. \\
Ισοδύναμα πρέπει $x > -1$ και $x < 1$. Οπότε η εξίσωση ορίζεται για $x \in (-1, 1)$.

\item Για $x \in (-1, 1)$ η εξίσωση γράφεται ισοδύναμα 
\begin{align*}
\log(x + 1) &= \log\left(\frac{1}{2}\right) - \log(1 - x) \Leftrightarrow \\
\log(x + 1) + \log(1 - x) &= \log(1) - \log(2) \Leftrightarrow \\
\log[(x + 1)(1 - x)] &= \log\left(\frac{1}{2}\right) \Leftrightarrow \\
(x + 1)(1 - x) &= \frac{1}{2} \Leftrightarrow \\
1 - x^2 &= \frac{1}{2} \Leftrightarrow \\
-x^2 &= \frac{-1}{2} \Leftrightarrow \\
x^2 &= \frac{1}{2} \Leftrightarrow \\
x &= \pm\frac{1}{\sqrt{2}} = \pm\frac{\sqrt{2}}{2} \in (-1, 1), \text{ δεκτές και οι δύο λύσεις.}
\end{align*}
\end{alist}
